%----------------------------------------------------------------------------------------
%	PLANTILLA DE INFORMES DE LABORATORIO
%	UNGS
%	LABORATORIO DE FÍSICA
%%%%%%%%%%%%%%%%%%%%%%%%%%%%%%%%%%%%%%%%%%%%%%%%%%%%%%

%%%%%%%%%%%%%%%%%%%%%%%%%%%%%%%%%%%%%%%%%%%%%%%%
% COMANDOS GENERALES
% En general estos comandos no se modifican al menos que querramos cambiar el aspecto del paper
%%%%%%%%%%%%%%%%%%%%%%%%%%%%%%%%%%%%%%%%%%%%%%%%
%Tipo de documento que estamos escribiendo:
\documentclass{article}  
% Paquetes de uso general:
%\usepackage[scaled=1]{helvet}
\usepackage{amsmath}
\usepackage[spanish,es-tabla]{babel}
\usepackage{graphicx}
\usepackage[utf8x]{inputenc}
\usepackage[T1]{fontenc}
\usepackage{multicol}
\usepackage{blindtext}
%Definición de márgenes:
\usepackage{vmargin}
\setpapersize{A4}
\setmargins{2.5cm}  % margen izquierdo
{1.0cm}             % margen superior
{16.5cm}            % anchura del texto
{23.42cm}           % altura del texto
{10pt}              % altura de los encabezados
{1cm}               % espacio entre el texto y los encabezados
{0pt}               % altura del pie de página
{1.0cm}             % espacio entre el texto y el pie de página

% Formato APA para la bibliografía
\usepackage[square,numbers]{natbib} 

%Otros comandos:
%\setlength\parindent{0pt} % Removes all indentation from paragraphs
\setlength{\parindent}{2em}
\setlength{\parskip}{1em}
\renewcommand{\baselinestretch}{1.25}

%%%%%%%%%%%%%%%%%%%%%%%%%%%%%%%%%%%%%%%%%%%%%%%%%%
% DOCUMENTO PRINCIPAL
%%%%%%%%%%%%%%%%%%%%%%%%%%%%%%%%%%%%%%%%%%%%%%%%%%
% ENCABEZADO
%%%%%%%%%%%%%%%%%%%%%%%%%%%%%%%%%%%%%%%%%%%%%%%%%%%
\begin{document}
\begin{center}
\includegraphics[scale=0.20]{LogoNuevoUNGS}\\
\vspace{0.5cm}
%%%%%% TITULO %%%%%%%%%%%%%%%%%%%%%%%%%%%%%%%%%%
%% El título del trabajo debe ser específico e informativo, y en lo posible agudo y provocador. Con él debemos dar una idea clara del tema estudiado.
\LARGE{Plantilla para informes de Laboratorio Avanzado}\\
\vspace{0.5cm}
%%%%%%%%%%%%%%%%%%%%%%%%%%%%%%%%%%%%%%%%%%%%%%%%
%%%%% AUTORES %%%%%%%%%%%%%%%%%%%%%%%%
%% MODIFICAR NOMRBRES Y DIRECCIONES DE E-MAIL
\small{A. Mengano$^1$, B. Fulano$^2$ y C. Montoto$^3$
\\
$^1$mengano@campus.ungs.edu.ar, 			  		     $^2$fulano@campus.ungs.edu.ar,         	 	      $^3$montoto@campus.ungs.edu.ar
}\\
\vspace{0.5cm}
\normalsize{Universidad Nacional de General Sarmiento \\ Laboratorio Avanzado\\}
\vspace{0.5cm}

%%%%%%%%%%%%%%%%%%%%%%%%%%%%%%%%%%%%
{Fecha de entrega: \today}\\
\vspace{0.5cm}
\end{center}
%%%%%%%%%%%%%%%%%%%%%%%%%%%%%%%%%%%%%
%%% TEXTO PRINCIPAL   %%%%%%%%%%%%%%
%%%%%%%%%%%%%%%%%%%%%%%%%%%%%%%%%%%%%
%%% RESUMEN %%%%%%%%%%%%%%%%%%%%%%%%%
\begin{abstract}
\begin{center}
\begin{minipage}{\dimexpr\paperwidth-10cm}

El resumen del informe debe dar un adelanto de lo que se leerá en el cuerpo del mismo, en lo posible en no más de 100 palabras. Aquí debemos indicar con concisión el tema del trabajo, referirnos sucintamente a la metodología seguida y destacar los resultados más importantes obtenidos. Los consejos para la redacción de informes fueron sutilmente robados de E. Rodríguez, profesor de la casa.


\end{minipage}
\end{center}
\end{abstract}

\bigskip

%%%%%%%%%%%%%%%%%%%%%%%%%%%%%%%%%%%%%%%%%%%%%%%%%%%%%%%%%%%%%%%%%%% EMPIEZA DOS COLUMNAS
%\begin{multicols}{2}
%\vspace{-1.2cm}
%%%%%%%%%%%%%%%%%%%%%%%%%%%%%%%%%%%%%%%%%%%%%%%%%%%%%%%%%%%%%%%%%%%



%%%%%%%%%%%%%%%% INTRODUCCIÓN  %%%%%%%%%%%%%%%%%%%%%%%%%%%%%%%%%%%%%%
\section*{Introducción}

El informe de laboratorio es una acabada prueba de que hicimos un experimento, lo analizamos y comprendimos. Cuando redactamos el informe es cuando terminamos de ordenar nuestros datos, gráficos, anotaciones y sobre todo nuestras ideas. Debe ofrecer a los lectores un recuento claro y completo de las actividades experimentales realizadas, nuestras conclusiones y reflexiones de lo que hicimos. El informe debe ser, ante todo claro, y en lo posible breve. Debemos redactarlo en lenguaje preciso y ameno, tratando de atraer y retener la atención de los lectores. Hagamos el siguiente ejercicio: Son las doce de la noche y el
lector de nuestro informe tiene también como opciones hojear el diario o ver televisión. Nuestro trabajo entrará en competencia con estas alternativas solo si está cuidadosamente redactado y si en él expresamos nuestras ideas con claridad y concisión. Esto podemos lograrlo usando construcciones cortas y cuidando que las descripciones no den lugar a interpretaciones ambiguas, de manera que el lector no se vea obligado a tener que volver sobre lo leído. Recordemos que no estaremos al lado de nuestro lector para hacerle aclaraciones a sus dudas y decirle que “donde escribimos una cosa”, en realidad, “quisimos decir otra”.

El informe no debe ser considerado como un documento que se presenta con el solo fin para que el profesor juzgue el trabajo realizado, sino que debe ser pensado como un texto que sea capaz de mostrar que hemos ganado la habilidad de comunicar por escrito nuestras ideas y resultados. Con esto en mente, los informes que se realizan en los cursos básicos de laboratorio son un muy buen entrenamiento para mejorar nuestra redacción y con ella nuestra capacidad de comunicar temas científicos y técnicos. Aquí damos algunas pautas y sugerencias sobre cómo organizar un informe de laboratorio.  

\textbf{En esta sección de introducción} debemos orientar al lector hacia el tema de estudio y la motivación por hacerlo elegido. Para esto es aconsejable que incluyamos un marco teórico–experimental del tema que estudiamos, con referencias adecuadas (ver Referencias) que lleven rápidamente a los antecedentes del problema y que destaquen la conexión de esas ideas con el trabajo realizado. Estas referencias deben orientar al lector hacia el \textit{estado del arte} del tema. Asimismo debemos enunciar claramente el propósito u objetivo del experimento. Quizás necesite incluir ecuaciones, lo cual se realiza de acuerdo al modelo siguiente, ver ecuación (\ref{ecuacion1}). Para escribir ecuaciones en línea o nombres de variables, puede hacerlo ensanguchando la expersión por signo pesos, por ejemplo $A$, $B$, $A+B=C$, etc.

%%%%%%%%%%%%%%%% ECUACION %%%%%%%%%%%%%%%%%%%%
\begin{equation}
\label{ecuacion1}
x(t)=x_{0}+v_{0}\: t+\frac{1}{2}\: a\: t^{2}
\end{equation}
%%%%%%%%%%%%%%%%%%%%%%%%%%%%%%%%%%%%%%%%%%%%%%

%%%%%%%%%%%%%%%% EXPERIMENTAL  %%%%%%%%%%%%%%%%%%%%%%%%%%%%%%%%%%%%%%
\section*{Método experimental}

También se puede llamar Materiales y Métodos. En esta sección describimos los procedimientos seguidos y el instrumental usado. Es útil incluir un esquema del diseño experimental elegido. Para esto puede recurrirse a diagramas esquemáticos que muestren las características más importantes del arreglo experimental y la disposición relativa de los instrumentos. Es una buena práctica indicar también cuáles variables se miden directamente, cuáles se obtienen indirectamente y a cuáles tomamos como datos de otras fuentes (parámetros físicos, constantes, etc.). También es aconsejable describir las virtudes y limitaciones del diseño experimental, analizar las fuentes de errores e individualizar las que aparezcan como las más críticas.



%%%%%%%%%%%%%%%% RESULTADOS  %%%%%%%%%%%%%%%%%%%%%%%%%%%%%%%%%%%%%%
\section*{Resultados}

Los resultados deben presentarse preferiblemente en forma de gráficos. En lo posible evitemos la inclusión de tablas de datos, a menos que sean sustanciales. Los datos del experimento deben estar diferenciados de otros datos que puedan incluirse para comparación y tomados de otras fuentes. Como práctica invariante, debemos expresar resultados con sus incertidumbres, en lo posible especificando cómo las calculamos. 

Para incluir tablas básicas, utilice los comandos de más abajo --- ver tabla \ref{tabla}, por ejemplo. Para modificar la tabla puede hacerlo en el modelo de abajo o recurrir a páginas que generan el código de una tabla diseñada previamente (ver tutorial).

%%%%%%%%%%% TABLA %%%%%%%%%%%%%%%%%%%%%%%%%%%%%
\begin{table}
\centering
\caption{\label{tabla}Un ejemplo de tabla.}
\bigskip
\begin{tabular}{l|r}
Ítem & Cantidad \\\hline
Chocolates & 42 \\
Caramelos & 13
\end{tabular}
\end{table}
%%%%%%%%%%%%%%%%%%%%%%%%%%%%%%%%%%%%%%%%%%%%%%%%%

Acá también hay un ejemplo de cómo incluir figuras. Si ven el archivo (main.tex) debajo de este párrafo van a encontrar el código general. Deben modificar el nombre de archivo, el nombre de la figura y el pie de figura. Para citarla en el texto se hará de la siguiente manera, ver figura \ref{potencia1}.


%%%%%%%%%%%%%%% FIGURA %%%%%%%%%%%%%%%%%%%%%%%%%%%%%%%
\begin{figure}
\centering
\includegraphics[width=0.7\textwidth]{potencia}
\caption{\label{potencia1} Un ejemplo de pie de figura.}
\end{figure}
%%%%%%%%%%%%%%%%%%%%%%%%%%%%%%%%%%%%%%%%%%%%%%%%%%%%%%%

%%%%%%%%%%%%%%%% DISCUSIÓN  %%%%%%%%%%%%%%%%%%%%%%%%%%%%%%%%%%%%%%
\section*{Análisis y Discusión}
 
Se puede agrupar esta sección con la anterior, en una sola que se titularía ``Resultados y Discusiones''. En esta parte debemos explicitar el análisis de los datos obtenidos. Aquí se analizan, por ejemplo, las dependencias observadas entre las variables, la comparación de los datos con un modelo propuesto, o las similitudes y discrepancias observadas con otros resultados. Si el trabajo además propone un modelo que trate de dar cuenta de los datos obtenidos, es decir, si el modelo es original del trabajo, su descripción debe quedar lo más clara posible; o bien, si se usó un modelo tomado de otros trabajos, debe citarse la fuente consultada. Si fuera necesaria una comparación de nuestros resultados con otros resultados previos, resaltemos similitudes y diferencias de los materiales, métodos y procedimientos empleados, para así poner en mejor contexto tal comparación.


%%%%%%%%%%%%%%%% CONCLUSIÓN  %%%%%%%%%%%%%%%%%%%%%%%%%%%%%%%%%%%%%%
\section*{Conclusiones}

En esta sección tenemos que comentar objetivamente qué hemos aprendido del experimento realizado, y sintetizar las consecuencias e implicancias que encontramos asociadas a nuestros resultados. Podemos decir que un buen informe es aquel que demuestra el mayor número de conclusiones (correctas) alcanzadas a partir de los datos obtenidos. 

%%%%%%%%%%%%%%%% APENDICE  %%%%%%%%%%%%%%%%%%%%%%%%%%%%%%%%%%%%%%
\section*{Apéndice}

Algunas veces son necesarios para la mejor comprensión de alguna parte del informe. Por lo general no es conveniente distraer al lector con muchos cálculos, despejes de términos y propagaciones de errores en la mitad del texto, así que este lugar puede ser propicio para estas consideraciones. En el texto principal deberemos orientar al lector para que consulte estos apéndices.

Algunos tips de escritura, para poner itálicas \textit{se hace así}. Para poner negrita \textbf{se hace así}. Para poner comillas ``se hace así''. Para poner listas se puede hacer así:

\begin{itemize}
  \item Una cosa.
  \item Otra cosa.
\end{itemize}


%%%%%%%%%%%%%%%% REFERENCIAS  %%%%%%%%%%%%%%%%%%%%%%%%%%%%%%%%%%%%%%
\section*{Referencias}

Las referencias bibliográficas se ordenan al final del informe. Deben contener el nombre de los autores de las publicaciones (artículos en revistas \cite{bb} o libros \cite{aa}) citados en el texto, el título de los trabajos; el nombre de la revista o editorial que los publicó; además se debe incluir los datos que ayuden a la identificación de los mismos: volumen donde están incluidos, capítulo, página, fecha de publicación, etc \cite{bb,aa,cc}. Se citan en el texto como se muestra en este ejemplo y se completa la cita en la parte de abajo de este texto. En esta sección, las referencias pueden ser sin numerar (escritas en forma genérica al estilo de una bibliografía):

\bigskip
\noindent Feynmann, R. (1965). \textit{Introduction to Physics (2da edición)}. Montevideo: editorial cuchuflito.

\bigskip
\noindent Sharp, T. (2004). The Physics of Physics. \textit{Journal of Physics}, 44(8), 19-23.

\bigskip
\noindent o también pueden ser numeradas, como en este caso (las cuales se citan en el texto):

%%%%% REFERENCIAS CITADAS %%%%%%%%%%%%%%%%%%
\begingroup
\renewcommand{\section}[2]{}

\begin{thebibliography}{}

\bibitem{bb} Sharp, T. (2004). The Physics of Physics. \textit{Journal of Physics}, 44(8), 19-23.

\bibitem{aa} Feynmann, R. (1965). \textit{Introduction to Physics (2da edición)}. Montevideo: editorial cuchuflito.

\bibitem{cc} Anónimo (1996). \textit{APA style citations of electronic sources}. Recuperado desde: \texttt{http://www.cas.usf.edu/}.

\end{thebibliography}
\endgroup



%% LIBRO: Autor. (Año de publicación.) Título en itálicas (edición). Lugar de publicación: Casa publicadora.

%%Revistas académicas (Journals): Autor o autores. (Año de publicación). Título del artículo. Título de la revista en itálicas, volumen(número de ejemplar), páginas.

%Web site: Autor. (año). Título del documento o del sitio en itálicas. Recuperado de lugar o dirección electrónica. 


%%%%%%%%%%%%%%%%%%%%%%%%%%%%%%%%%%%%%%%%%%%%%%%%%%%%%%%%%%%%%%%%%%%%
%TERMINAR dos columnas
%\end{multicols}
        
\end{document}